\SeventhArticle{5}

Since the formula $\Delta(AR+\allowbreak A'R'+\allowbreak A'R'+\allowbreak A''R''+\allowbreak \text{etc.}+\allowbreak A^{(n)}R^{(n)})$ gives the integral $\int y\,\mathrm{d}x$ from $x=g$ to $x=g+\Delta$, or the integral $\Delta\int y\,\mathrm{d}t$ from $t=0$ to $t=1$, exactly whenever $y$, expanded in a series, does not go beyond the power $t^n$, but only approximately whenever $y$ goes further, we still have to show how to determine the error introduced by the immediately following terms. Let us denote in general by $k^{(m)}$ the difference between the true value of the integral $\int t^m\,\mathrm{d}t$ from $t=0$ to $t=1$ and the value yielded by the formula, so that
\[
\begin{aligned}
k&=1-R-R'-R''-R'''-\text{etc.}-R^{(n)}\\
k'&=\frac12-\frac1n(R'+2R''+3R'''+\text{etc.}+nR^{(n)})\\
k''&=\frac13-\frac1{nn}(R'+4R''+9R'''+\text{etc.}+nnR^{(n)})\\
k'''&=\frac14-\frac1{n^2}(R'+8R''+27R'''+\text{etc.}\ n^3R^{(n)})
\end{aligned}
\]
etc. It is therefore clear that if $y$ is expanded in the series
\[
K+K't+K''tt+K'''t^3+\text{etc.}
\]
the difference between the true value of the integral $\int y\,\mathrm{d}t$ and the approximate value supplied by the formula is expressed by
\[
Kk+K'k'+K''k''+K'''k'''+\text{etc.}
\]
But $k$, $k'$, $k''$, etc., up to $k^{(n)}$, evidently vanish automatically ($=0$); the correction to the approximate formula will therefore be
\[
K^{(n+1)}k^{(n+1)}+K^{(n+2)}k^{(n+2)}+K^{(n+3)}k^{(n+3)}+\text{etc.}
\]

We shall investigate the nature of the quantities $k^{(n+1}$, $k^{(n+2)}$, etc., more closely below. Here it is enough to give the numerical values of the first or the second for each value of $n$, so that the degree of accuracy supplied by the approximate formula can be estimated from them.
\SourceLeaf{179}

\begin{longtable}{@{}l@{\quad}l@{}}
$\text{For }n=1\ \text{we have }$ & $k''=-\frac16,\quad k'''=-\frac14,\quad k''''=-\frac3{10}$ \\[16pt]
$\text{For }n=2\ \text{we find }$ & $k'''=0,\quad k''''=-\frac1{120},\quad k^{\mathrm{V}}=-\frac1{48}$ \\[16pt]
$\text{For }n=3\ \text{one obtains }$ & $k''''=-\frac1{270},\quad k^{\mathrm{V}}=-\frac1{108}$ \\[16pt]
$\text{For }n=4\ldotp\ldotp\ldotp\ $ & $k^{\mathrm{V}}=0,\quad k^{\mathrm{VI}}=-\frac1{2688},\quad k^{\mathrm{VII}}=-\frac1{768}$ \\[16pt]
$\text{For }n=5\ldotp\ldotp\ldotp\ $ & $k^{\mathrm{VI}}=-\frac{11}{52500},\quad k^{\mathrm{VII}}=-\frac{11}{15000}$ \\[16pt]
$\text{For }n=6\ldotp\ldotp\ldotp\ $ & $k^{\mathrm{VII}}=0,\quad k^{\mathrm{VIII}}=-\frac1{38880},\quad k^{\mathrm{IX}}=-\frac1{8640}$ \\[16pt]
$\text{For }n=7\ldotp\ldotp\ldotp\ $ & $k^{\mathrm{VIII}}=-\frac{167}{10588410},\quad k^{\mathrm{IX}}=-\frac{167}{2352980}$ \\[16pt]
$\text{For }n=8\ldotp\ldotp\ldotp\ $ & $\begin{aligned}[t]&k^{\mathrm{IX}}=0,\quad k^{\mathrm{X}}=-\frac{37}{17301504},\\[3pt]&k^{\mathrm{XI}}=-\frac{37}{3145728}\end{aligned}$ \\[16pt]\noalign{\vskip 8pt}
$\text{For }n=9\ldotp\ldotp\ldotp\ $ & $\begin{aligned}[t]&k^{\mathrm{X}}=-\frac{865}{631351908},\\[3pt]&k^{\mathrm{XI}}=-\frac{865}{114791256}\end{aligned}$ \\[16pt]\noalign{\vskip 8pt}
$\text{For }n=10\ldotp\ldotp\ldotp\ $ & $\begin{aligned}[t]&k^{\mathrm{XI}}=0,\quad k^{\mathrm{XII}}=-\frac{26927}{136500000000},\\[3pt]&k^{\mathrm{XIII}}=-\frac{26927}{21000000000}\end{aligned}$ \\[16pt]
\end{longtable}




Here, for an even value of $n$, we observe in every case that $k^{(n+1)}=0$, and also $k^{(n+3)}=\frac{n+3}{2}k^{(n+2)}$; for an odd value of $n$, however, one obtains in every case $k^{(n+2)}=\frac{n+2}{2}k^{(n+1)}$. The reason for these results is easily drawn from the following considerations.

Let us denote in general by $l^{(m)}$ the difference between the true value of the integral $\int(t-\frac12)^m\,\mathrm{d}t$ from $t=0$ to $t=1$ and the value yielded by the approximate formula, so that
\[
\begin{aligned}
l^{(m)}={}&\int(t-\tfrac12)^m\,\mathrm{d}t
-\bigl[(-\tfrac12)^mR+(\tfrac1n-\tfrac12)^mR'
\\ &\qquad +(\tfrac2n-\tfrac12)^mR''+(\tfrac3n-\tfrac12)^mR'''+\text{etc.}\\
&\qquad +(\tfrac12-\tfrac1n)^mR^{(n-1)}+(\tfrac12)^mR^{(n)}\bigr]
\end{aligned}
\]
the integral being taken from $t=0$ to $t=1$. Evidently, for an odd value of $m$, both the true and the approximate values of the integral vanish. Thus $l'=0$, $l'''=0$, $l^{\mathrm{V}}=0$, $l^{\mathrm{VII}}=0$, etc.; or, in general, $l^{(m)}=0$ for an odd value of $m$. For an even value of $m$, however, we give the formula the following form
\[
\begin{aligned}
l^{(m)}={}&\frac1{2^m(m+1)}-\frac2{n^m}\bigl((\tfrac12n)^mR
+(\tfrac12n-1)^mR'\\ &\qquad +(\tfrac12n-2)^mR''+\text{etc.}\\
&\qquad +2^mR^{(\frac12n-2)}+R^{(\frac12n-1)}\bigr)
\end{aligned}
\]
if $n$ is also even; or the following form
\[
\begin{aligned}
l^{(m)}={}&\frac1{2^m}\bigl(\frac1{m+1}-\frac2{n^m}\bigl(n^mR
+(n-2)^mR'\\ &\qquad +(n-4)^mR''+\text{etc.}\\
&\qquad +3^mR^{(\frac12n-\frac32)}+R^{(\frac12n-\frac12)}\bigr)\bigr)
\end{aligned}
\]
if $n$ is also odd.
\par\SourceLeaf{180}
Thus, if the expansion of $y$ in a series in powers of $t-\frac12$ gives
\[
y=L+L'(t-\tfrac12)+L''(t-\tfrac12)^2+L'''(t-\tfrac12)^3+\text{etc.}
\]
the correction to be applied to the approximate value of the integral $\int y\,\mathrm{d}t$ from $t=0$ to $t=1$ will be
\[
Ll+L''l''+L''''l''''+L^{\mathrm{VI}}l^{\mathrm{VI}}+\text{etc.}
\]
\Needspace{14\baselineskip}or rather, since $l^{(m)}$ necessarily vanishes for any integer value of $m$ not greater than $n$, the correction will be
\[
L^{(n+2)}l^{(n+2)}+L^{(n+4)}l^{(n+4)}+L^{(n+6)}l^{(n+6)}+\text{etc.}
\]
for even $n$, or
\[
L^{(n+1)}l^{(n+1)}+L^{(n+3)}l^{(n+3)}+L^{(n+5)}l^{(n+5)}+\text{etc.}
\]
for odd $n$.

The corrections $l^{(m)}$ are now very easily expressed in terms of $k^{(m)}$, and conversely. For, since
\[
(t-\tfrac12)^m=t^m-\tfrac12m\mathbin{.}t^{m-1}
+\tfrac14\mathbin{.}\frac{m(m-1)}{1\mathbin{.}2}t^{m-2}+\text{etc.}
\]
one obtains
\[
l^{(m)}=k^{(m)}-\tfrac12m k^{(m-1)}
+\tfrac14\mathbin{.}\frac{m(m-1)}{1\mathbin{.}2}k^{(m-2)}+\text{etc.}
\]
Likewise one obtains
\[
k^{(m)}=l^{(m)}+\tfrac12m l^{(m-1)}
+\tfrac14\mathbin{.}\frac{m(m-1)}{1\mathbin{.}2}l^{(m-1)}+\text{etc.}
\]



From the latter formula, the terms in which $l$ has an odd index will be removed; both formulas, however, need be continued only as far as the index $n+1$, inclusive. We shall therefore evidently have
\[
\begin{aligned}
&\textit{for $n$ even}\\
&\qquad k^{(n+1)}=0\\
&\qquad k^{(n+2)}=l^{(n+2)}\\
&\qquad k^{(n+3)}=\frac{n+3}{2}\mathbin{.}l^{(n+2)}\\
&\textit{for $n$ odd}\\
&\qquad k^{(n+1)}=l^{(n+1)}\\
&\qquad k^{(n+2)}=\frac{n+2}{2}\mathbin{.}l^{(n+1)}
\end{aligned}
\]
from which follow the observations stated above.
\SourceLeaf{181}
\SeventhArticle{6}

Generally speaking, therefore, when applying the method of Cotes, it will be preferable to give $n$ an even value, that is, to use an odd number of terms. Accuracy increases very little if we advance from an even value of $n$ to the next higher odd value, since the error remains of the same order, though its coefficient is somewhat smaller. Conversely, in passing from an odd value of $n$ to the next even value, the order of the error increases by two, and a more considerable reduction of the coefficient increases the accuracy further. Thus, if five terms are used, or for $n=4$, the error is approximately expressed by $-\frac{1}{2688}K^6$ or by $-\frac{1}{2688}L^6$. If $n=5$ is set, the error will be approximately $-\frac{11}{52500}K^6$ or $-\frac{11}{52500}L^6$, and is thus not even reduced to half the former error. By contrast, putting $n=6$ gives for the error approximately $=-\frac{1}{38880}K^8$ or $=-\frac{1}{38880}L^8$; the accuracy is increased all the more, the faster the series in which the function is expanded already converges on its own.

\SeventhArticle{7}

After these preliminary remarks on the method of Cotes, we proceed to a general investigation, abandoning the condition that the values of $x$ form an arithmetic progression. We therefore undertake the problem of determining the value of the integral $\int y\,\mathrm{d}x$ between given limits from a number of given values of $y$, either exactly or as closely as possible. Let us suppose that the integral is to be taken from $x=g$ to $x=g+\Delta$, and let us introduce, in place of $x$, another variable $t=\frac{x-g}{\Delta}$, so that we must investigate the integral $\Delta\int y\,\mathrm{d}t$ from $t=0$ to $t=1$. Let the $n+1$ given values of $y$, namely $A,\ A',\ A'',\ A'''\ldotp\ldotp\ldotp\ldotp A^{(n)}$, correspond to the distinct values of $t$, namely $a,\ a',\ a'',\ a'''\ldotp\ldotp\ldotp\ldotp a^{(n)}$; and let us denote by $Y$ the following polynomial of degree $n$:
\[
\begin{aligned}
& A\;\frac{(t-a')(t-a'')(t-a''')\ldotp\ldotp\ldotp\ldotp(t-a^{(n)})}
{(a-a')(a-a'')(a-a''')\ldotp\ldotp\ldotp\ldotp(a-a^{(n)})}\\[6pt]
&{}+A'\;\frac{(t-a)(t-a'')(t-a''')\ldotp\ldotp\ldotp\ldotp(t-a^{(n)})}
{(a'-a)(a'-a'')(a'-a''')\ldotp\ldotp\ldotp\ldotp(a'-a^{(n)})}\\[6pt]
&{}+A''\;\frac{(t-a)(t-a')(t-a''')\ldotp\ldotp\ldotp\ldotp(t-a^{(n)})}
{(a''-a)(a''-a')(a''-a''')\ldotp\ldotp\ldotp\ldotp(a''-a^{(n)})}\\[6pt]
&{}+\text{etc.}\\[6pt]
&{}+A^{(n)}\;\frac{(t-a)(t-a')(t-a'')\ldotp\ldotp\ldotp\ldotp(t-a^{(n-1)})}
{(a^{(n)}-a)(a^{(n)}-a')(a^{(n)}-a'')\ldotp\ldotp\ldotp\ldotp(a^{(n)}-a^{(n-1)})}
\end{aligned}
\]
The values of this function, when $t$ is set equal to one of the quantities $a,\allowbreak \ a',\allowbreak \ a'',\allowbreak \ a'''\ldotp\ldotp\ldotp\ldotp a^{(m)}$, evidently coincide with the corresponding values of the function $y$. Hence,\SourceLeaf{182} just as in Article 2, we conclude that $Y$ is identical with $y$ whenever $y$ is also a polynomial of degree not exceeding $n$, or that it can at least take the place of $y$ if $y$, when expanded in a series in powers of $t$, converges so rapidly that terms of higher orders may be neglected.


